%______________________________________________________________________
%   Notes to authors
% Please disable your editor's auto newline wrapping feature.
% Please format \itemize sections
%______________________________________________________________________

\chapter{Post-Processing Scripts} \label{Chapter:scripts}



%______________________________________________________________________
\section{parseVarTrackerVars.py}

The Python scripts described in this chapter parses the output from a simulation
with the \verb|VarTracker| debugging option enabled.
\verb|VarTracker| dumps the value of specific variables, for specific
tasks and materials, at every timestep, into the \verb|sus|
 output file. These scripts parse that raw log file into per-timestep,
per-task, per-variable data files, and then into plots.

The scripts live in \verb|src/scripts/VarTracker| and are used in a
two-stage pipeline:

\begin{enumerate}
  \item \verb|parseVarTrackerVars.py| parses the raw \verb|sus| output into a
    directory of small, clean data files.
  \item \verb|plotVarTrackerVars.py| turns those data files into X-Y
    plots, one plot per (timestep, variable), with one line per task.  These plots can be hand-tuned after the fact.
\end{enumerate}

\subsection{Purpose}
\verb|parseVarTrackerVars.py| was written to aid the developer by showing  
how a variable changes during a timestep. This is different than examining the variables that are in the \verb|uda| since multiple tasks can modify a variable during a timestep

\subsection{Usage}
\begin{Verbatim}[fontsize=\footnotesize]
  parseVarTrackerVars.py sus_output [-t TASK] [-v VARIABLE] [-m MATERIAL]
\end{Verbatim}

\begin{itemize}
  \item \verb|sus_output| -- the \verb|sus| log file to parse (required,
    positional).
  \item \verb|-t|/\verb|--task| -- a task name to search for. Repeatable.
    If omitted entirely, defaults to all tasks.
  \item \verb|-v|/\verb|--variable| -- a variable name to search for.
    Repeatable. Default: all variables.
  \item \verb|-m|/\verb|--material| -- a material index to search for.
    Repeatable. Default: all materials.
\end{itemize}

\noindent Supplying any \verb|-t|/\verb|-v|/\verb|-m| flag at all replaces that
flag's default list entirely (the values given are not merged with the
defaults). For example, to track only \verb|press_CC| for material 0:

\begin{Verbatim}[fontsize=\footnotesize]
  parseVarTrackerVars.py mySimulation.log -v press_CC -m 0
\end{Verbatim}

\subsection{Output}
All output is written under a \verb|parsedVarTracker/| directory created
in the current working directory:

\begin{Verbatim}[fontsize=\footnotesize]
  parsedVarTracker/<timestep>/<task>/<variable>_<material>
\end{Verbatim}

\noindent Each such file is plain text, safe to read directly or to feed to
\verb|gnuplot| (comment lines begin with \verb|#|, which \verb|gnuplot|
skips natively):

\begin{Verbatim}[fontsize=\footnotesize]
  # Variable: sp_vol_CC, DW 1, Patch 0, Matl 1
  # Time: 0
  # x,  y,  z,  value
   16, 10, 0, 8.510598642398e-01
   16, 11, 0, 8.513706469716e-01
   ...
\end{Verbatim}

\noindent The \verb|# Time:| line gives the physical simulation time for that
timestep (parsed from the log's \verb|Timestep N  Time=...| header),
distinct from the timestep index used in the file's own directory path.
Intermediate index files the script builds while parsing are deleted
before it exits -- only the final per-(timestep, task, variable, material)
files remain.

%______________________________________________________________________
\section{plotVarTrackerVars.py} \label{Section:plotVarTrackerVars}

\subsection{Purpose}
\verb|plotVarTrackerVars.py| produces the plots (one PNG file per timestep/variable, one line per
task) via \verb|gnuplot|, and leaves behind the \verb|gnuplot| scripts used to generate them
so the plots' appearance can be tuned by hand afterward, without re-running the Python script.

\subsection{Usage}
\begin{Verbatim}[fontsize=\footnotesize]
  plotVarTrackerVars.py [parsed_dir] [-o OUTPUT]
\end{Verbatim}
 This script depends on \verb|gnuplot| to be in your path; no Python plotting package is needed.

\subsection{Output}
\begin{Verbatim}[fontsize=\footnotesize]
  <output>/<variable>/template.gp
  <output>/<variable>/regenerate_plots.sh
  <output>/<variable>/<timestep>.png
\end{Verbatim}

Each variable gets its own subdirectory containing three kinds of files:

\begin{itemize}
  \item \verb|template.gp| -- one shared \verb|gnuplot| template for
    \emph{every} timestep's plot of this variable (terminal, font, key
    placement, grid, etc.). This is the file to hand-edit: any styling
    change here is picked up by every plot of this variable the next time
    \verb|regenerate_plots.sh| is run. Other variables' templates are
    completely independent -- editing one variable's \verb|template.gp|
    never affects another's.
  %
  \item \verb|regenerate_plots.sh| -- a generated, executable shell script
    that re-renders every PDF for this variable from \verb|template.gp|.
    It defines \verb|parsed_dir| and \verb|out_dir| shell variables once
    at the top (so the whole \verb|parsedVarTracker| tree can be moved
    and only those two lines need updating), then loops over one
    \verb|gnuplot -e "..."| invocation per timestep, each setting the
    output file, plot title, axis labels, and plot command for that one
    timestep before applying \verb|template.gp|.
  %
  \item \verb|<timestep>.png| -- the rendered plot itself.
\end{itemize}

\noindent To change the appearance of every plot of one variable: edit that
variable's \verb|template.gp|, then run:
\begin{Verbatim}[fontsize=\footnotesize]
  bash <output>/<variable>/regenerate_plots.sh
\end{Verbatim}
This re-renders only that variable's plots; no Python involved, and no
effect on any other variable.

\noindent The plot title includes the variable name and, when available, the
physical simulation time read from the \verb|# Time:| comment
\verb|parseVarTrackerVars.py| writes into each data file, for example:
\begin{Verbatim}[fontsize=\footnotesize]
  sp_vol_CC -- Timestep 2 (Time=3e-08)
\end{Verbatim}

%______________________________________________________________________

\section{Example}

\noindent Suppose you the following debugging \verb|VarTracker| section was in the \verb|ups| file and a simulation was run.

\begin{Verbatim}[fontsize=\scriptsize]

   <Scheduler>
      <VarTracker>
        <output formatFlag="scientific" precision="12"/>
        
        <start_time> 0 </start_time>
        <end_time>  1 </end_time>
        <start_index> [50,-1,50]  </start_index>
        <end_index>   [50,100,50] </end_index>
  
        <var label="press_CC"        matl_id="0" dw= "NewDW" />
        <var label="grad_P_YFC"      matl_id="0" dw= "NewDW" singleTask="ICE::computeVel_FC"/>
        <var label="sp_vol_CC"       matl_id="0" dw= "NewDW"/>
        <var label="temp_CC"         matl_id="0" dw= "NewDW" singleTask="ICE::conservedtoPrimitive_Vars:afterAdvection"/>
        
        <locations after_exec = "true"/>
        <!--<locations before_exec = "true"/> -->
        
        <task name="ICE::computeEquilibrationPressure"/>
        <task name="ICE::computeVel_FC"/>
         <task name="ICE::conservedtoPrimitive_Vars:afterAdvection"/>
      </VarTracker>
    </Scheduler>   
\end{Verbatim}

After the simulation was finished the user would then run 
\begin{Verbatim}[fontsize=\footnotesize]
  src/scripts/VarTracker/parseVarTrackerVars.py sus_out.txt
\end{Verbatim}

to parse \verb|sus_out.txt| to produce a series of files:

\begin{Verbatim}[fontsize=\scriptsize]
parsedVarTracker/
|-- 0
|   |-- ICE::computeVel_FC
|   |   |-- grad_P_YFC_0
|   |   `-- sp_vol_CC_0
|   `-- ICE::conservedtoPrimitive_Vars:afterAdvection
|       |-- press_CC_0
|       |-- sp_vol_CC_0
|       `-- temp_CC_0
|-- 1
|   |-- ICE::computeVel_FC
|   |   |-- grad_P_YFC_0
|   |   `-- sp_vol_CC_0
|   `-- ICE::conservedtoPrimitive_Vars:afterAdvection
|       |-- press_CC_0
|       |-- sp_vol_CC_0
|       `-- temp_CC_0

 <snip>
\end{Verbatim}

\noindent To plot these the variables the user would then run:
\begin{Verbatim}[fontsize=\footnotesize]
  src/scripts/VarTracker/plotVarTrackerVars.py
\end{Verbatim}

\noindent to produce a series of png files, 1 per timestep per variable, see Fig.~\ref{fig:varTrackerPlot}.

\begin{Verbatim}[fontsize=\scriptsize]
parsedVarTracker/plots/
|-- grad_P_YFC
|   |-- 0.png
|   |-- 1.png
|   |-- 2.png
|   |-- 3.png
|   |-- regenerate_plots.sh
|   `-- template.gp         << This can be hand tuned
|-- press_CC
|   |-- 0.png
|   |-- 1.png
|   |-- 2.png
|   |-- 3.png
|   |-- regenerate_plots.sh
|   `-- template.gp         << This can be hand tuned
|-- sp_vol_CC
|   |-- 0.png
|   |-- 1.png
|   |-- 2.png
|   |-- 3.png
|   |-- regenerate_plots.sh
|   `-- template.gp         << This can be hand tuned
`-- temp_CC
    |-- 0.png
    |-- 1.png
    |-- 2.png
    |-- 3.png
    |-- regenerate_plots.sh
    `-- template.gp         << This can be hand tuned
\end{Verbatim}

%
\begin{figure}[h]
  \centering
  \fbox{ \includegraphics[scale=0.5]{varTracker.png} }
  \caption{ Example plot after the VarTracker output has been post processed. }
  \label{fig:varTrackerPlot}
\end{figure}
